\section{Haftung und Rechtsverfolgung im Netz}

\subsection{Ausgangslage und Cyberrisiken}

\begin{example}
    \textbf{Cyberrisiken (z.B. Ransomware-Angriffe):} \\
    Angreifer verschlüsseln geschäftskritische Daten oder legen IT-Systeme lahm und fordern Lösegeld für die Freigabe (oder drohen bei \textit{Double Extortion} mit der Veröffentlichung gestohlener Daten). Dies führt zu massiven Betriebsunterbrüchen, Wiederherstellungskosten, Reputationsschäden und möglichen Ansprüchen geschädigter Dritter.
\end{example}

\begin{important}
    \textbf{Grundproblem der Rechtsverfolgung im Netz:} \\
    Die eigentlichen Täter (z.B. Hacker oder Urheber von Rechtsverletzungen) agieren meist \textbf{anonym} und aus dem \textbf{Ausland}. Sie sind rechtlich und faktisch kaum greifbar. \\
    $\rightarrow$ Geschädigte und Behörden nehmen daher für die Haftung und Rechtsdurchsetzung primär greifbare Akteure in der Schweiz ins Visier, insbesondere betroffene \textbf{Unternehmen, deren Organe sowie IT-Dienstleister und Provider}.
\end{important}

\subsection{Grundvoraussetzungen der Haftung}

\begin{theorem}
    \textbf{Prinzip der gesetzlichen Grundlage:} \\
    Jede rechtliche Verantwortlichkeit (straf- wie zivilrechtlich) setzt eine \textbf{gesetzliche Grundlage} voraus. Für das Greifen einer Rechtsfolge müssen grundsätzlich drei Ebenen erfüllt sein:
    \begin{enumerate}
        \item \textbf{Objektiver Tatbestand:} Äussere Merkmale einer Rechtsverletzung oder Straftat (Tathandlung, Taterfolg, Kausalzusammenhang).
        \item \textbf{Subjektiver Tatbestand (Verschulden):} Innere Einstellung bzw. Vorwerfbarkeit:
              \begin{itemize}
                  \item \textit{Vorsatz:} Wissen und Wollen der Tatbestandsverwirklichung.
                  \item \textit{Eventualvorsatz:} Ernstnehmen der Möglichkeit des Erfolgs und Inkaufnahme desselben ("billigend in Kauf nehmen").
                  \item \textit{Fahrlässigkeit:} Pflichtwidrige Unvorsichtigkeit (Missachtung der nach den Umständen und persönlichen Verhältnissen gebotenen Sorgfalt).
              \end{itemize}
        \item \textbf{Fehlen von Rechtfertigungsgründen:} Keine gesetzliche Erlaubnis, keine gültige Einwilligung des Verletzten, keine Notwehr oder Notstand.
    \end{enumerate}
\end{theorem}

\subsubsection{Strafrechtliche Verantwortlichkeit}

\begin{remark}
    \textbf{Typische Cyber-Straftatbestände der Angreifer (StGB):} \\
    Im Strafrecht gilt ein strenges Legalitätsprinzip (\textit{nulla poena sine lege}, Art. 1 StGB); strafbar ist in der Regel nur vorsätzliches (inkl. eventualvorsätzliches) Handeln, sofern das Gesetz Fahrlässigkeit nicht ausdrücklich unter Strafe stellt. Angreifer machen sich insbesondere strafbar wegen:
    \begin{itemize}
        \item \textbf{Unbefugtes Eindringen in ein Datenverarbeitungssystem (Hacking, Art. 143bis StGB):} Eindringen in fremde, gegen unbefugten Zugriff besonders gesicherte Systeme auf elektronischem Weg.
        \item \textbf{Unbefugte Datenbeschaffung (Art. 143 StGB):} Beschaffen elektronisch gespeicherter, nicht für den Täter bestimmter und besonders gesicherter Daten in Bereicherungsabsicht.
        \item \textbf{Datenbeschädigung (Art. 144bis StGB):} Unbefugtes Ändern, Löschen oder Unbrauchbarmachen elektronisch gespeicherter Daten (z.B. Verschlüsselung durch Ransomware, DDoS-Attacken) sowie das Herstellen/Verbreiten von Schadsoftware (Malware).
        \item \textbf{Erpressung (Art. 156 StGB) und Nötigung (Art. 181 StGB):} Erzwingen einer Lösegeldzahlung (z.B. in Kryptowährungen) durch Androhung ernstlicher Nachteile (Datenverlust oder Veröffentlichung sensibler Daten).
        \item \textbf{Betrügerischer Missbrauch einer Datenverarbeitungsanlage (Art. 147 StGB):} Vermögensverschiebung durch unrichtige, unvollständige oder unbefugte Verwendung von Daten (z.B. Phishing).
    \end{itemize}
\end{remark}

\subsubsection{Zivilrechtliche Haftung (Schadenersatz)}

\begin{definition}
    \textbf{Voraussetzungen der zivilrechtlichen Haftung (Art. 41 ff. OR / Art. 97 ff. OR):} \\
    Damit ein Geschädigter zivilrechtlich Schadenersatz fordern kann, müssen \textbf{vier Voraussetzungen kumulativ} nachgewiesen werden:
    \begin{enumerate}
        \item \textbf{Vermögensschaden:} Nachgewiesene, unfreiwillige Vermögensverminderung (Differenztheorie: Bestandesabnahme, Passivenzunahme oder entgangener Gewinn).
        \item \textbf{Widerrechtlichkeit oder Vertragswidrigkeit:}
              \begin{itemize}
                  \item \textit{Ausservertraglich (Delikt, Art. 41 OR):} Verletzung eines absolut geschützten Rechtsguts (z.B. Eigentum, Persönlichkeit, Urheberrecht) oder Verstoss gegen eine spezifische Schutznorm.
                  \item \textit{Vertraglich (Art. 97 OR):} Verletzung einer Pflicht aus einem bestehenden Schuldverhältnis (z.B. Hosting-Vertrag, SLA, Sorgfaltspflichten).
              \end{itemize}
        \item \textbf{Adäquater Kausalzusammenhang:} Das schädigende Verhalten muss nach dem gewöhnlichen Lauf der Dinge und der allgemeinen Lebenserfahrung geeignet sein, einen Erfolg von der Art des eingetretenen herbeizuführen.
        \item \textbf{Verschulden (Absicht oder Fahrlässigkeit):} Subjektive Vorwerfbarkeit (im Vertragsrecht wird das Verschulden gesetzlich vermutet, d.h. der Schuldner muss sich exkulpieren, Art. 97 Abs. 1 OR; im Deliktsrecht muss der Geschädigte das Verschulden beweisen).
    \end{enumerate}
\end{definition}

\subsection{Haftungsverteilung, Rückgriff und Organverantwortung}

\subsubsection{Haftung im Unternehmen und Regress}

\begin{remark}
    \textbf{Haftung für eigenes Verhalten, Angestellte und Hilfspersonen:}
    \begin{itemize}
        \item \textbf{Primäre Unternehmenshaftung:} Unternehmen haften im Aussenverhältnis gegenüber Kunden und Dritten primär selbst für Schäden, die durch ihre Organe (Art. 55 ZGB / Art. 722 OR), ihre Arbeitnehmer oder beigezogene Hilfspersonen (z.B. externe IT-Dienstleister, Subunternehmer) verursacht werden.
        \item \textbf{Vertragliche Hilfspersonenhaftung (Art. 101 OR):} Wer die Erfüllung einer Schuldpflicht durch eine Hilfsperson (z.B. Mitarbeiter, Cloud-Anbieter) vornehmen lässt, haftet dem Vertragspartner für den von der Hilfsperson in Ausübung ihrer Verrichtung verursachten Schaden ohne eigenes Verschulden.
        \item \textbf{Deliktische Geschäftsherrenhaftung (Art. 55 OR):} Als Geschäftsherr haftet das Unternehmen für Schäden durch Arbeitnehmer/Hilfspersonen, sofern es nicht nachweist, dass es alle nach den Umständen gebotene Sorgfalt angewendet hat (\textbf{Cura in eligendo, instruendo et custodiendo}).
        \item \textbf{Eingeschränkter Rückgriff (Regress auf Arbeitnehmer, Art. 321e OR):} Hat das Unternehmen den Schaden im Aussenverhältnis gedeckt, ist der interne Rückgriff auf fehlbare Angestellte (z.B. Mitarbeiter klickt auf Phishing-Mail) stark eingeschränkt. Bei leichter Fahrlässigkeit (Alltagsversehen) trägt in der Regel der Arbeitgeber das Betriebsrisiko; ein voller oder teilweiser Rückgriff ist meist nur bei grober Fahrlässigkeit oder Vorsatz unter Berücksichtigung des Berufsrisikos, des Bildungsgrades und der Lohnhöhe zulässig.
    \end{itemize}
\end{remark}

\subsubsection{Verantwortung des Verwaltungsrats und Sorgfaltsmassstab}

\begin{definition}
    \textbf{Verantwortung des Verwaltungsrats (Art. 716a \& Art. 754 OR):} \\
    Der Verwaltungsrat (VR) trägt die unübertragbare und unentziehbare Oberleitung des Unternehmens, einschliesslich der Ausgestaltung des Risikomanagements und der Internen Kontrolle (IKS). \\
    Delegiert der VR operative Aufgaben (z.B. IT-Security an die Geschäftsleitung oder den CISO), haftet er für die sorgfältige:
    \begin{enumerate}
        \item \textbf{Auswahl (\textit{Cura in eligendo}):} Einsetzen fachlich und persönlich qualifizierter Personen.
        \item \textbf{Instruktion (\textit{Cura in instruendo}):} Klare Weisungen, Richtlinien, Sicherheitskonzepte und Zielvorgaben.
        \item \textbf{Überwachung (\textit{Cura in custodiendo}):} Regelmässige Kontrolle, Audits, Reporting und Einschreiten bei Missständen.
    \end{enumerate}
\end{definition}

\begin{definition}
    \textbf{Three Lines of Defence Modell (Risikomanagement \& Compliance):} \\
    Anerkanntes Governance-Modell zur systematischen Steuerung von Risiken (insb. Cyber- und IT-Risiken) und zur Sicherstellung der Compliance im Unternehmen:
    \begin{enumerate}
        \item \textbf{1st Line of Defence (Operatives Management \& Fachabteilungen):} Die operativen Einheiten und der IT-Betrieb identifizieren, bewerten und steuern Risiken direkt im Tagesgeschäft durch technische und organisatorische Schutzmassnahmen (z.B. Patch-Management, Firewalls, Zugriffskontrollen).
        \item \textbf{2nd Line of Defence (Risikomanagement, Compliance \& IT-Security/CISO):} Überwacht und unterstützt die erste Linie, gibt Sicherheitsrichtlinien und Rahmenwerke vor und prüft die Einhaltung gesetzlicher sowie interner Vorgaben.
        \item \textbf{3rd Line of Defence (Interne Revision / Internal Audit):} Prüft völlig unabhängig von den operativen Prozessen die Wirksamkeit der ersten und zweiten Verteidigungslinie und rapportiert direkt an den Verwaltungsrat (bzw. das Audit Committee).
    \end{enumerate}
\end{definition}

\begin{important}
    \textbf{Sorgfaltsmassstab und Branchenstandards (Ausschluss der Fahrlässigkeit):} \\
    Es gibt im Netz keine absolute (100\%-ige) Sicherheit gegen Cyberangriffe. \\
    \textbf{Grundsatz:} Wer anerkannte \textbf{Branchenstandards} (Stand der Technik, z.B. ISO 27001, NIST Cybersecurity Framework, IKT-Minimalstandard des Bundes) angemessen umsetzt, eine funktionierende Governance (z.B. \textit{Three Lines of Defence}) pflegt und die gebotene Sorgfalt walten lässt, \textbf{haftet nicht wegen Fahrlässigkeit}, selbst wenn es dennoch zu einem erfolgreichen Angriff oder Sicherheitsvorfall kommt.
\end{important}

\subsection{Provider-Haftung und Rechtsdurchsetzung}

\begin{theorem}
    \textbf{Grundprinzip der Provider-Verantwortlichkeit:} \\
    In der Schweiz existiert kein allgemeines Spezialgesetz für sämtliche Online-Provider (anders als z.B. der Digital Services Act in der EU); es gelten die allgemeinen zivil- und strafrechtlichen Regeln (mit einzelnen Spezialnormen wie Art. 39d URG). \\
    Als Faustregel gilt: \textbf{Je näher ein Provider an den Daten und Inhalten ist (Kontroll- und Einflussmöglichkeit), desto weitreichender sind seine Prüf-, Sorgfalts- und Handlungspflichten.}
\end{theorem}

\begin{remark}
    \textbf{Abstufung nach Provider-Typen:}
    \begin{itemize}
        \item \textbf{Access-Provider (Zugangsvermittler, z.B. Swisscom, Sunrise):} \\
              Vermitteln lediglich den technischen Zugang zum Internet (reine Durchleitung / \textit{Mere Conduit}) und haben weder Kenntnis noch Kontrolle über die transportierten Datenpakete. \\
              $\rightarrow$ Sie haften \textbf{in der Regel nicht} für Rechtsverletzungen Dritter, die über ihren Zugang begangen werden, und sind grundsätzlich auch nicht zu Netzsperren wegen zivilrechtlicher Ansprüche verpflichtet (Bundesgerichtspraxis; Ausnahmen nur bei expliziter gesetzlicher Grundlage wie Geldspiel-Sperrlisten nach BGS).
        \item \textbf{Hosting-Provider (Speicherplatzanbieter \& Plattformen, z.B. Webhoster, Cloud-Anbieter, Social Media):} \\
              Speichern fremde Inhalte im Auftrag ihrer Nutzer auf eigenen Servern und halten diese zum Abruf bereit. \\
              $\rightarrow$ Sie haben zwar \textbf{keine allgemeine Vorab-Prüfpflicht} (keine Pflicht zur proaktiven Überwachung sämtlicher hochgeladener Inhalte), sind jedoch zur \textbf{Entfernung rechtswidriger Inhalte verpflichtet}, sobald sie konkrete Kenntnis von einer Rechtsverletzung erlangen (\textbf{Notice-and-Takedown}-Prinzip, vgl. auch Art. 39d URG für Urheberrechte mit \textit{Stay-down}-Pflicht bei besonderer Gefahr). \\
              $\rightarrow$ Für Unterlassungs- und Beseitigungsansprüche (z.B. Löschung persönlichkeitsverletzender Beiträge gemäss Art. 28 ZGB) können Hosting-Provider als Mitwirkende direkt ins Recht gefasst werden, auch ohne eigenes Verschulden.
        \item \textbf{Content-Provider (Inhalteanbieter):} \\
              Stellen eigene Inhalte ins Netz oder machen sich fremde Inhalte zu eigen. \\
              $\rightarrow$ Sie tragen die \textbf{volle eigene zivil- und strafrechtliche Verantwortlichkeit} für die publizierten Inhalte.
    \end{itemize}
\end{remark}